
@@M0@@

L'homomorphisme @@M1@@ se déduit de l'homomorphisme : @@M2@@ et coïncide avec celui qui se déduit de @@M3@@; l'homomorphisme : @@M4@@ se déduit du carré commutatif (1). La propriété a) est donc évidente.

L'homomorphisme : @@M5@@ considéré dans b) peut être défini comme suit: soit @@M6@@, soit @@M7@@ un élément de @@M8@@ remontant @@M9@@ (il en existe car @@M10@@). On associe à @@M11@@ l'image de @@M12@@ par l'homomorphisme composé : @@M13@@. Par ailleurs, soit @@M14@@ l'idéal d'augmentation de : @@M15@@ ; le carré commutatif (2) donne un homomorphisme : @@M16@@, d'où un homomorphisme : @@M17@@. Or @@M18@@ ; par suite, il suffit de vérifier que les homomorphismes construits sont inverses l'un de l'autre, ce qui se fait sans difficulté.

\textbf{Définition 2.3.} \uline{Soit} @@M19@@ \uline{un morphisme lissifiable}. \uline{Le complexe} @@M20@@, \uline{défini à isomorphisme canonique près dans la catégorie dérivée} @@M21@@, \uline{est appelé} complexe cotangent relatif \uline{de} @@M22@@ \uline{sur} @@M23@@ (\uline{ou de} @@M24@@).

\textbf{Proposition 2.4.} \uline{Soit} @@M25@@ \uline{un morphisme d'intersection complète} (1.1) \uline{lissifiable}. \uline{Alors} @@M26@@ \uline{est parfait}. \uline{Plus précisément}, \uline{pour toute}
